
\documentclass{article}
\usepackage{colortbl}
\usepackage{makecell}
\usepackage{multirow}
\usepackage{supertabular}

\begin{document}

\newcounter{utterance}

\centering \large Interaction Transcript for game `function\_detective', experiment `numbers\_passive\_examples', episode 37 with Qwen3.6{-}35B{-}A3B{-}FP8{-}without{-}reasoning.
\vspace{24pt}

{ \footnotesize  \setcounter{utterance}{1}
\setlength{\tabcolsep}{0pt}
\begin{supertabular}{c@{$\;$}|p{.15\linewidth}@{}p{.15\linewidth}p{.15\linewidth}p{.15\linewidth}p{.15\linewidth}p{.15\linewidth}}
    \# & \multicolumn{2}{c}{Player} && \multicolumn{2}{c}{Game Master} \\
    \hline

    \theutterance \stepcounter{utterance}  
    & & & \multicolumn{4}{p{0.6\linewidth}}{
        \cellcolor[rgb]{0.9,0.9,0.9}{
            \makecell[{{p{\linewidth}}}]{
                \texttt{\tiny{[P1$\langle$GM]}}
                \texttt{I have a hidden Python function for you to figure out.} \\
\\ 
\texttt{Your job is to infer what it does and then submit a correct Python implementation.} \\
\\ 
\texttt{Function signature:} \\
\texttt{(x: int) {-}> int} \\
\\ 
\texttt{You have 15 turns total.} \\
\\ 
\texttt{I control which valid input{-}output examples you see. You can request them one at a time.} \\
\\ 
\texttt{You must make your best SOLVE attempt no later than the second{-}to{-}last turn.} \\
\texttt{Do not wait until the very end.} \\
\\ 
\texttt{You can use these actions:} \\
\texttt{{-} JSON format described at the end of the prompt} \\
\texttt{{-} SOLVE: ```python ... ```} \\
\\ 
\texttt{Do not add import statements in your guess.} \\
\texttt{Anything you might need is already available, so just write the function body in hypothesis and return the answer.} \\
\\ 
\texttt{Preloaded examples:} \\
\texttt{1. (1000, 505)} \\
\\ 
\texttt{Scoring:} \\
\texttt{{-} correctness is required for any non{-}zero score} \\
\texttt{{-} in interactive modes, solving with fewer information requests is better} \\
\\ 
\texttt{Input formatting reminders:} \\
\texttt{{-} strings must be quoted} \\
\texttt{{-} booleans must be True/False} \\
\texttt{{-} use commas between multiple values} \\
\texttt{{-} do not write any extra action tags in the same message} \\
\\ 
\texttt{Format example:} \\
\texttt{You write:} \\
\texttt{\{'current\_mode': <probing/confirming/don't know; choose whichever best describes your current state in the task>, 'input\_rationale': <brief description of why you chose this hypothesis>, 'hypothesis': <either the exact string "unknown" if you have no guess yet OR your current best guess as runnable Python matching the given signature, with a one{-}line docstring stating what it does. Never write placeholder code, if you cannot commit to a guess, use "unknown". Match the given signature.} \\
\texttt{e.g.} \\
\texttt{def f(x):} \\
\texttt{	"""Returns the absolute value of x."""} \\
\texttt{	return x if x >= 0 else {-}x} \\
\\ 
\texttt{Use "unknown" if you have no guess.>, 'confidence\_score': <score between 1{-}100. 100 means highest confidence in your hypothesis>\}} \\
\\ 
\texttt{Output only one valid json per turn.} \\
            }
        }
    }
     \\ \\

    \theutterance \stepcounter{utterance}  
    & \multicolumn{4}{p{0.6\linewidth}}{
        \cellcolor[rgb]{0.561,0.737,0.561}{
            \makecell[{{p{\linewidth}}}]{
                \texttt{\tiny{[P1$\rangle$GM]}}
                \texttt{\{'current\_mode': 'probing', 'input\_rationale': 'I have one example: f(1000) = 505. This looks like it might be related to summing digits or some arithmetic property. 1000 has digits 1,0,0,0. Sum is 1. 505 is not obviously related. Let me try a few more inputs to narrow down the function. I will try small integers first to see if there is a pattern like f(x) = x + something, or f(x) = reverse(x), or sum of digits, etc.', 'hypothesis': 'unknown', 'confidence\_score': 10\}} \\
            }
        }
    }
    & & \\ \\

    \theutterance \stepcounter{utterance}  
    & & & \multicolumn{4}{p{0.6\linewidth}}{
        \cellcolor[rgb]{0.9,0.9,0.9}{
            \makecell[{{p{\linewidth}}}]{
                \texttt{\tiny{[P1$\langle$GM]}}
                \texttt{Game Violation: ParseError: Response is not valid JSON.} \\
\texttt{Output ONLY:} \\
\texttt{One JSON with confidence\_score, current\_mode, input\_rationale, and hypothesis keys ...} \\
\texttt{OR} \\
\texttt{SOLVE: ```python ... ```} \\
            }
        }
    }
     \\ \\

    \theutterance \stepcounter{utterance}  
    & \multicolumn{4}{p{0.6\linewidth}}{
        \cellcolor[rgb]{0.561,0.737,0.561}{
            \makecell[{{p{\linewidth}}}]{
                \texttt{\tiny{[P1$\rangle$GM]}}
                \texttt{\{'current\_mode': 'probing', 'input\_rationale': 'I have one example: f(1000) = 505. This looks like it might be related to summing digits or some arithmetic property. 1000 has digits 1,0,0,0. Sum is 1. 505 is not obviously related. Let me try a few more inputs to narrow down the function. I will try small integers first to see if there is a pattern like f(x) = x + something, or f(x) = reverse(x), or sum of digits, etc.', 'hypothesis': 'unknown', 'confidence\_score': 10\}} \\
            }
        }
    }
    & & \\ \\

    \theutterance \stepcounter{utterance}  
    & & & \multicolumn{4}{p{0.6\linewidth}}{
        \cellcolor[rgb]{0.9,0.9,0.9}{
            \makecell[{{p{\linewidth}}}]{
                \texttt{\tiny{[P1$\langle$GM]}}
                \texttt{Game Violation: ParseError: Response is not valid JSON.} \\
\texttt{Output ONLY:} \\
\texttt{One JSON with confidence\_score, current\_mode, input\_rationale, and hypothesis keys ...} \\
\texttt{OR} \\
\texttt{SOLVE: ```python ... ```} \\
            }
        }
    }
     \\ \\

    \theutterance \stepcounter{utterance}  
    & \multicolumn{4}{p{0.6\linewidth}}{
        \cellcolor[rgb]{0.561,0.737,0.561}{
            \makecell[{{p{\linewidth}}}]{
                \texttt{\tiny{[P1$\rangle$GM]}}
                \texttt{\{'current\_mode': 'probing', 'input\_rationale': 'I have one example: f(1000) = 505. I need more data points to identify the pattern. I will request f(1), f(2), and f(10) to see if there is a simple arithmetic or digit{-}based relationship.', 'hypothesis': 'unknown', 'confidence\_score': 10\}} \\
            }
        }
    }
    & & \\ \\

    \theutterance \stepcounter{utterance}  
    & & & \multicolumn{4}{p{0.6\linewidth}}{
        \cellcolor[rgb]{0.9,0.9,0.9}{
            \makecell[{{p{\linewidth}}}]{
                \texttt{\tiny{[P1$\langle$GM]}}
                \texttt{Game Violation: ParseError: Response is not valid JSON.} \\
\texttt{Output ONLY:} \\
\texttt{One JSON with confidence\_score, current\_mode, input\_rationale, and hypothesis keys ...} \\
\texttt{OR} \\
\texttt{SOLVE: ```python ... ```} \\
            }
        }
    }
     \\ \\

    \theutterance \stepcounter{utterance}  
    & \multicolumn{4}{p{0.6\linewidth}}{
        \cellcolor[rgb]{0.561,0.737,0.561}{
            \makecell[{{p{\linewidth}}}]{
                \texttt{\tiny{[P1$\rangle$GM]}}
                \texttt{\{"current\_mode": "probing", "input\_rationale": "I have one example: f(1000) = 505. I need more data points to identify the pattern. I will request f(1), f(2), and f(10) to see if there is a simple arithmetic or digit{-}based relationship.", "hypothesis": "unknown", "confidence\_score": 10\}} \\
            }
        }
    }
    & & \\ \\

    \theutterance \stepcounter{utterance}  
    & & & \multicolumn{4}{p{0.6\linewidth}}{
        \cellcolor[rgb]{0.9,0.9,0.9}{
            \makecell[{{p{\linewidth}}}]{
                \texttt{\tiny{[P1$\langle$GM]}}
                \texttt{OUTPUT: (1, 5)} \\
            }
        }
    }
     \\ \\

    \theutterance \stepcounter{utterance}  
    & \multicolumn{4}{p{0.6\linewidth}}{
        \cellcolor[rgb]{0.561,0.737,0.561}{
            \makecell[{{p{\linewidth}}}]{
                \texttt{\tiny{[P1$\rangle$GM]}}
                \texttt{\{"current\_mode": "probing", "input\_rationale": "I have two examples: f(1) = 5 and f(1000) = 505. Let's analyze these. f(1)=5 is interesting. f(1000)=505. Let's try to see if there's a pattern. Maybe it's related to the sum of digits? sum\_digits(1)=1, sum\_digits(1000)=1. Not obvious. Maybe it's x + something? 1+4=5, 1000{-}495=505. Not consistent. Let's try another input, like f(2) or f(10).", "hypothesis": "unknown", "confidence\_score": 10\}} \\
            }
        }
    }
    & & \\ \\

    \theutterance \stepcounter{utterance}  
    & & & \multicolumn{4}{p{0.6\linewidth}}{
        \cellcolor[rgb]{0.9,0.9,0.9}{
            \makecell[{{p{\linewidth}}}]{
                \texttt{\tiny{[P1$\langle$GM]}}
                \texttt{OUTPUT: (99, 54)} \\
            }
        }
    }
     \\ \\

    \theutterance \stepcounter{utterance}  
    & \multicolumn{4}{p{0.6\linewidth}}{
        \cellcolor[rgb]{0.561,0.737,0.561}{
            \makecell[{{p{\linewidth}}}]{
                \texttt{\tiny{[P1$\rangle$GM]}}
                \texttt{\{"current\_mode": "probing", "input\_rationale": "I have three examples: f(1)=5, f(99)=54, f(1000)=505. Let's look at the differences: f(1){-}1=4, f(99){-}99={-}45, f(1000){-}1000={-}495. Not linear. Let's look at digits. 1 {-}> 5. 99 {-}> 54. 1000 {-}> 505. Maybe it's related to reversing digits? reverse(1)=1, reverse(99)=99, reverse(1000)=0001=1. No. Maybe it's sum of digits times something? sum(1)=1, 1*5=5. sum(99)=18, 18*3=54. sum(1000)=1, 1*505=505. Multipliers are 5, 3, 505. Not consistent. Let's try f(10) to get more data.", "hypothesis": "unknown", "confidence\_score": 10\}} \\
            }
        }
    }
    & & \\ \\

    \theutterance \stepcounter{utterance}  
    & & & \multicolumn{4}{p{0.6\linewidth}}{
        \cellcolor[rgb]{0.9,0.9,0.9}{
            \makecell[{{p{\linewidth}}}]{
                \texttt{\tiny{[P1$\langle$GM]}}
                \texttt{OUTPUT: ({-}2, 4)} \\
            }
        }
    }
     \\ \\

    \theutterance \stepcounter{utterance}  
    & \multicolumn{4}{p{0.6\linewidth}}{
        \cellcolor[rgb]{0.561,0.737,0.561}{
            \makecell[{{p{\linewidth}}}]{
                \texttt{\tiny{[P1$\rangle$GM]}}
                \texttt{\{"current\_mode": "probing", "input\_rationale": "I have four examples: f(1)=5, f(99)=54, f(1000)=505, f({-}2)=4. Let's look at the absolute values. f(1)=5, f(2)=4 (from f({-}2)=4, assuming symmetry or similar behavior for magnitude). f(99)=54, f(1000)=505. Let's check if it's related to the sum of digits of the absolute value. sum\_digits(1)=1 {-}> 5. sum\_digits(2)=2 {-}> 4. sum\_digits(99)=18 {-}> 54. sum\_digits(1000)=1 {-}> 505. The multipliers are 5, 2, 3, 505. Not consistent. Let's try f(10) to see if there's a pattern with powers of 10.", "hypothesis": "unknown", "confidence\_score": 10\}} \\
            }
        }
    }
    & & \\ \\

    \theutterance \stepcounter{utterance}  
    & & & \multicolumn{4}{p{0.6\linewidth}}{
        \cellcolor[rgb]{0.9,0.9,0.9}{
            \makecell[{{p{\linewidth}}}]{
                \texttt{\tiny{[P1$\langle$GM]}}
                \texttt{OUTPUT: (0, 5)} \\
            }
        }
    }
     \\ \\

    \theutterance \stepcounter{utterance}  
    & \multicolumn{4}{p{0.6\linewidth}}{
        \cellcolor[rgb]{0.561,0.737,0.561}{
            \makecell[{{p{\linewidth}}}]{
                \texttt{\tiny{[P1$\rangle$GM]}}
                \texttt{\{"current\_mode": "probing", "input\_rationale": "I have five examples: f(1)=5, f(99)=54, f(1000)=505, f({-}2)=4, f(0)=5. Let's look at the relationship between x and f(x). f(0)=5. f(1)=5. f({-}2)=4. f(99)=54. f(1000)=505. Let's check if f(x) = sum of digits of x + 4? sum(0)+4=4!=5. sum(1)+4=5. sum(99)+4=22!=54. No. Let's check if f(x) is related to the number of digits or something else. Let's try f(10) and f(11) to get more data on small numbers.", "hypothesis": "unknown", "confidence\_score": 10\}} \\
            }
        }
    }
    & & \\ \\

    \theutterance \stepcounter{utterance}  
    & & & \multicolumn{4}{p{0.6\linewidth}}{
        \cellcolor[rgb]{0.9,0.9,0.9}{
            \makecell[{{p{\linewidth}}}]{
                \texttt{\tiny{[P1$\langle$GM]}}
                \texttt{OUTPUT: ({-}1, 4)} \\
            }
        }
    }
     \\ \\

    \theutterance \stepcounter{utterance}  
    & \multicolumn{4}{p{0.6\linewidth}}{
        \cellcolor[rgb]{0.561,0.737,0.561}{
            \makecell[{{p{\linewidth}}}]{
                \texttt{\tiny{[P1$\rangle$GM]}}
                \texttt{\{"current\_mode": "probing", "input\_rationale": "I have six examples: f(1)=5, f(99)=54, f(1000)=505, f({-}2)=4, f(0)=5, f({-}1)=4. Let's look at the pattern more closely. f(0)=5, f(1)=5. f({-}1)=4, f({-}2)=4. It seems like for x >= 0, f(x) might be related to something, and for x < 0, it's different. Let's check f(2) and f(10) to see if the pattern for positive numbers holds. Also, let's check if f(x) = sum of digits of x + 4 for x >= 0? sum(0)+4=4!=5. sum(1)+4=5. sum(99)+4=22!=54. No. Let's try f(2) and f(10).", "hypothesis": "unknown", "confidence\_score": 10\}} \\
            }
        }
    }
    & & \\ \\

    \theutterance \stepcounter{utterance}  
    & & & \multicolumn{4}{p{0.6\linewidth}}{
        \cellcolor[rgb]{0.9,0.9,0.9}{
            \makecell[{{p{\linewidth}}}]{
                \texttt{\tiny{[P1$\langle$GM]}}
                \texttt{OUTPUT: ({-}1000, {-}495)} \\
            }
        }
    }
     \\ \\

    \theutterance \stepcounter{utterance}  
    & \multicolumn{4}{p{0.6\linewidth}}{
        \cellcolor[rgb]{0.561,0.737,0.561}{
            \makecell[{{p{\linewidth}}}]{
                \texttt{\tiny{[P1$\rangle$GM]}}
                \texttt{\{"current\_mode": "probing", "input\_rationale": "I have seven examples: f(1)=5, f(99)=54, f(1000)=505, f({-}2)=4, f(0)=5, f({-}1)=4, f({-}1000)={-}495. Let's analyze the relationship between x and f(x). For positive x: f(1)=5, f(99)=54, f(1000)=505. For negative x: f({-}1)=4, f({-}2)=4, f({-}1000)={-}495. For zero: f(0)=5. Let's look at the absolute values. |f(1)|=5, |f(99)|=54, |f(1000)|=505. |f({-}1)|=4, |f({-}2)|=4, |f({-}1000)|=495. f(0)=5. Let's check if f(x) = sum of digits of x + 4 for x >= 0? sum(0)+4=4!=5. sum(1)+4=5. sum(99)+4=22!=54. No. Let's check if f(x) = reverse(x) + something? reverse(1)=1,} \\
            }
        }
    }
    & & \\ \\

    \theutterance \stepcounter{utterance}  
    & & & \multicolumn{4}{p{0.6\linewidth}}{
        \cellcolor[rgb]{0.9,0.9,0.9}{
            \makecell[{{p{\linewidth}}}]{
                \texttt{\tiny{[P1$\langle$GM]}}
                \texttt{Game Violation: ParseError: Response is not valid JSON.} \\
\texttt{Output ONLY:} \\
\texttt{One JSON with confidence\_score, current\_mode, input\_rationale, and hypothesis keys ...} \\
\texttt{OR} \\
\texttt{SOLVE: ```python ... ```} \\
            }
        }
    }
     \\ \\

    \theutterance \stepcounter{utterance}  
    & \multicolumn{4}{p{0.6\linewidth}}{
        \cellcolor[rgb]{0.561,0.737,0.561}{
            \makecell[{{p{\linewidth}}}]{
                \texttt{\tiny{[P1$\rangle$GM]}}
                \texttt{\{"current\_mode": "probing", "input\_rationale": "I have seven examples: f(1)=5, f(99)=54, f(1000)=505, f({-}2)=4, f(0)=5, f({-}1)=4, f({-}1000)={-}495. Let's look at the relationship between x and f(x) more carefully. For positive x, f(x) seems to be related to the sum of digits. Let s(x) be the sum of digits of x. s(1)=1, f(1)=5. s(99)=18, f(99)=54. s(1000)=1, f(1000)=505. Notice that 5 = 1 * 5, 54 = 18 * 3, 505 = 1 * 505. The multipliers are 5, 3, 505. This doesn't seem consistent. Let's look at negative x. f({-}1)=4, f({-}2)=4, f({-}1000)={-}495. s(1)=1, f({-}1)=4. s(2)=2, f({-}2)=4. s(1000)=1, f({-}1000)={-}495. Let's} \\
            }
        }
    }
    & & \\ \\

    \theutterance \stepcounter{utterance}  
    & & & \multicolumn{4}{p{0.6\linewidth}}{
        \cellcolor[rgb]{0.9,0.9,0.9}{
            \makecell[{{p{\linewidth}}}]{
                \texttt{\tiny{[P1$\langle$GM]}}
                \texttt{Game Violation: ParseError: Response is not valid JSON.} \\
\texttt{Output ONLY:} \\
\texttt{One JSON with confidence\_score, current\_mode, input\_rationale, and hypothesis keys ...} \\
\texttt{OR} \\
\texttt{SOLVE: ```python ... ```} \\
            }
        }
    }
     \\ \\

    \theutterance \stepcounter{utterance}  
    & \multicolumn{4}{p{0.6\linewidth}}{
        \cellcolor[rgb]{0.561,0.737,0.561}{
            \makecell[{{p{\linewidth}}}]{
                \texttt{\tiny{[P1$\rangle$GM]}}
                \texttt{\{"current\_mode": "probing", "input\_rationale": "I have seven examples: f(1)=5, f(99)=54, f(1000)=505, f({-}2)=4, f(0)=5, f({-}1)=4, f({-}1000)={-}495. Let's look at the relationship between x and f(x) more carefully. For positive x, f(x) seems to be related to the sum of digits. Let s(x) be the sum of digits of x. s(1)=1, f(1)=5. s(99)=18, f(99)=54. s(1000)=1, f(1000)=505. Notice that 5 = 1 * 5, 54 = 18 * 3, 505 = 1 * 505. The multipliers are 5, 3, 505. This doesn't seem consistent. Let's look at negative x. f({-}1)=4, f({-}2)=4, f({-}1000)={-}495. s(1)=1, f({-}1)=4. s(2)=2, f({-}2)=4. s(1000)=1, f({-}1000)={-}495. Let's} \\
            }
        }
    }
    & & \\ \\

    \theutterance \stepcounter{utterance}  
    & & & \multicolumn{4}{p{0.6\linewidth}}{
        \cellcolor[rgb]{0.9,0.9,0.9}{
            \makecell[{{p{\linewidth}}}]{
                \texttt{\tiny{[P1$\langle$GM]}}
                \texttt{Game Violation: ParseError: Response is not valid JSON.} \\
\texttt{Output ONLY:} \\
\texttt{One JSON with confidence\_score, current\_mode, input\_rationale, and hypothesis keys ...} \\
\texttt{OR} \\
\texttt{SOLVE: ```python ... ```} \\
            }
        }
    }
     \\ \\

    \theutterance \stepcounter{utterance}  
    & \multicolumn{4}{p{0.6\linewidth}}{
        \cellcolor[rgb]{0.561,0.737,0.561}{
            \makecell[{{p{\linewidth}}}]{
                \texttt{\tiny{[P1$\rangle$GM]}}
                \texttt{\{"current\_mode": "probing", "input\_rationale": "I have seven examples: f(1)=5, f(99)=54, f(1000)=505, f({-}2)=4, f(0)=5, f({-}1)=4, f({-}1000)={-}495. Let's look at the relationship between x and f(x) more carefully. For positive x, f(x) seems to be related to the sum of digits. Let s(x) be the sum of digits of x. s(1)=1, f(1)=5. s(99)=18, f(99)=54. s(1000)=1, f(1000)=505. Notice that 5 = 1 * 5, 54 = 18 * 3, 505 = 1 * 505. The multipliers are 5, 3, 505. This doesn't seem consistent. Let's look at negative x. f({-}1)=4, f({-}2)=4, f({-}1000)={-}495. s(1)=1, f({-}1)=4. s(2)=2, f({-}2)=4. s(1000)=1, f({-}1000)={-}495. Let's} \\
            }
        }
    }
    & & \\ \\

    \theutterance \stepcounter{utterance}  
    & & & \multicolumn{4}{p{0.6\linewidth}}{
        \cellcolor[rgb]{0.9,0.9,0.9}{
            \makecell[{{p{\linewidth}}}]{
                \texttt{\tiny{[P1$\langle$GM]}}
                \texttt{Game Violation: ParseError: Response is not valid JSON.} \\
\texttt{Output ONLY:} \\
\texttt{One JSON with confidence\_score, current\_mode, input\_rationale, and hypothesis keys ...} \\
\texttt{OR} \\
\texttt{SOLVE: ```python ... ```} \\
            }
        }
    }
     \\ \\

    \theutterance \stepcounter{utterance}  
    & \multicolumn{4}{p{0.6\linewidth}}{
        \cellcolor[rgb]{0.561,0.737,0.561}{
            \makecell[{{p{\linewidth}}}]{
                \texttt{\tiny{[P1$\rangle$GM]}}
                \texttt{\{"current\_mode": "probing", "input\_rationale": "I have seven examples: f(1)=5, f(99)=54, f(1000)=505, f({-}2)=4, f(0)=5, f({-}1)=4, f({-}1000)={-}495. Let's look at the relationship between x and f(x) more carefully. For positive x, f(x) seems to be related to the sum of digits. Let s(x) be the sum of digits of x. s(1)=1, f(1)=5. s(99)=18, f(99)=54. s(1000)=1, f(1000)=505. Notice that 5 = 1 * 5, 54 = 18 * 3, 505 = 1 * 505. The multipliers are 5, 3, 505. This doesn't seem consistent. Let's look at negative x. f({-}1)=4, f({-}2)=4, f({-}1000)={-}495. s(1)=1, f({-}1)=4. s(2)=2, f({-}2)=4. s(1000)=1, f({-}1000)={-}495. Let's} \\
            }
        }
    }
    & & \\ \\

    \theutterance \stepcounter{utterance}  
    & & & \multicolumn{4}{p{0.6\linewidth}}{
        \cellcolor[rgb]{0.9,0.9,0.9}{
            \makecell[{{p{\linewidth}}}]{
                \texttt{\tiny{[P1$\langle$GM]}}
                \texttt{Game Violation: ParseError: Response is not valid JSON.} \\
\texttt{Output ONLY:} \\
\texttt{One JSON with confidence\_score, current\_mode, input\_rationale, and hypothesis keys ...} \\
\texttt{OR} \\
\texttt{SOLVE: ```python ... ```} \\
            }
        }
    }
     \\ \\

    \theutterance \stepcounter{utterance}  
    & \multicolumn{4}{p{0.6\linewidth}}{
        \cellcolor[rgb]{0.561,0.737,0.561}{
            \makecell[{{p{\linewidth}}}]{
                \texttt{\tiny{[P1$\rangle$GM]}}
                \texttt{\{"current\_mode": "probing", "input\_rationale": "I have seven examples: f(1)=5, f(99)=54, f(1000)=505, f({-}2)=4, f(0)=5, f({-}1)=4, f({-}1000)={-}495. Let's look at the relationship between x and f(x) more carefully. For positive x, f(x) seems to be related to the sum of digits. Let s(x) be the sum of digits of x. s(1)=1, f(1)=5. s(99)=18, f(99)=54. s(1000)=1, f(1000)=505. Notice that 5 = 1 * 5, 54 = 18 * 3, 505 = 1 * 505. The multipliers are 5, 3, 505. This doesn't seem consistent. Let's look at negative x. f({-}1)=4, f({-}2)=4, f({-}1000)={-}495. s(1)=1, f({-}1)=4. s(2)=2, f({-}2)=4. s(1000)=1, f({-}1000)={-}495. Let's} \\
            }
        }
    }
    & & \\ \\

    \theutterance \stepcounter{utterance}  
    & & & \multicolumn{2}{p{0.3\linewidth}}{
        \cellcolor[rgb]{0.9,0.9,0.9}{
            \makecell[{{p{\linewidth}}}]{
                \texttt{\tiny{[GM$|$GM]}}
                \texttt{Game Verdict: LOSS (turn limit reached)} \\
            }
        }
    }
    & & \\ \\

\end{supertabular}
}

\end{document}
